\section{Related Work}
\label{sec:related}

\paragraph{Distributed hash tables.} Chord \citep{stoica2001chord}, Pastry \citep{rowstron2001pastry}, and Kademlia \citep{maymounkov2002kademlia} established the structured-overlay model inherited by Freenet's ring. They organize peers by hashed identifiers in a ring or XOR-metric space and provide $O(\log n)$ lookup through maintained routing tables. Freenet stores mutable, contract-defined state at each key and constructs its topology through the mechanism in \cref{sec:routing}, without explicit finger tables.

\paragraph{Content-addressed storage.} The original Freenet \citep{clarke2001freenet} introduced content-hash keys (CHKs) for immutable data around 1999--2000. IPFS \citep{benet2014ipfs} later applied content addressing to web deployment over its Bitswap block exchange. In both systems, external naming or application layers provide mutability. Freenet contracts instead bind mutable state and its consistency rules to content-addressed code.

\paragraph{CRDTs.} Conflict-free Replicated Data Types \citep{shapiro2011crdt} provide the algebraic foundation for the contract layer. The contract interface implements the state-based (CvRDT) model while allowing applications to supply both the lattice and a summary/delta synchronization protocol. Antidote and AntidoteDB \citep{akkoorath2016cure} similarly provide CRDTs for geo-distributed databases, but expose a predefined CRDT vocabulary.

\paragraph{Bayou.} Bayou \citep{terry1995bayou} is the closest historical precedent for contracts. In 1995 it used application-supplied merge procedures on trusted replica servers to reconcile concurrent writes. Bayou procedures had database access and were configured per database. Freenet executes content-addressed merge code as sandboxed WebAssembly on untrusted hosts and restricts it to an idempotent commutative monoid. These restrictions allow replicas to converge without the platform inspecting the merge implementation or trusting its host.

\paragraph{Anti-entropy and delta sync.} Dynamo \citep{decandia2007dynamo} and its descendants use Merkle anti-entropy to reconcile replicas. Delta-state CRDTs \citep{almeida2018delta} formalize state-based synchronization that transfers only changes. Freenet follows the same general protocol shape, with contract-supplied lattices and delta encodings operating over a small-world routing layer instead of a fully connected replica set.

\paragraph{Local-first software and CRDT libraries.} Local-first software \citep{kleppmann2019local} favors convergence on user devices over server-mediated consistency. Automerge \citep{automerge} and Y.js \citep{yjs} provide composable CRDT structures such as JSON-like trees, sequences, and counters while leaving transport to the application. Those libraries guarantee convergence when applications remain within their primitives. Freenet lets each contract define its algebra and synchronization protocol, placing the corresponding correctness obligation on the contract author.

\paragraph{Secure Scuttlebutt.} SSB \citep{scuttlebutt} centers on per-user append-only signed logs replicated by gossip on demand. Such a feed can be represented as a Freenet contract whose merge is set union over signed records. Freenet generalizes the state structure and replication protocol beyond logs while targeting many of the same applications, including private messaging, blogging, and social feeds.

\paragraph{Hypercore.} The Hypercore protocol family \citep{hypercore} provides content-addressed, append-only signed logs replicated peer-to-peer, with Hyperbee and Hyperdrive adding higher-level data structures. A Hyperbee resembles a contract with library-defined merge semantics, while a local key-management agent over Hypercore serves a role similar to a delegate.

\paragraph{Small-world routing.} Kleinberg's result \citep{kleinberg2000small} provides the theoretical basis for the routing layer; the original Freenet \citep{clarke2001freenet} predated and informally anticipated parts of it. Symphony \citep{manku2003symphony} combines short-range ring links with $1/d$ long-range links sampled from a harmonic distribution. Freenet likewise protects its closest successor and predecessor, but continuously adjusts the long-range distribution: a joiner targets an under-represented region of its connections, and the receiver favors candidates that close large gaps. To our knowledge, this combination has not been analyzed in the structured-overlay literature.

\paragraph{Performance-aware routing.} Adaptive routing in peer-to-peer overlays has a long history. Freenet uses isotonic regression, a standard calibration method \citep{deleeuw2009isotone}, as a monotone distance-based baseline, then combines it in a weighted ensemble with Renegade's k-nearest-neighbor predictions \citep{renegade} to represent peer--contract interactions. The two components predict the same outcomes; Renegade is not trained as a residual correction to the isotonic model.

\paragraph{Blockchain platforms.} Ethereum \citep{wood2014ethereum} and successor smart-contract platforms execute a globally linearized replicated state machine. Freenet instead provides per-contract eventual consistency with application-defined merge. It is intended for interactive applications with mergeable updates and does not support applications that require a single global transaction order.

\paragraph{Federated systems.} Matrix \citep{matrix} is a server-based federated system with a similar application surface. Its state-resolution algorithm reconciles divergent room states by a deterministic application-defined rule. Freenet applies this approach in a peer-to-peer setting and permits arbitrary idempotent commutative monoids.

\paragraph{The original Freenet.} The original Freenet \citep{clarke2001freenet} provided anonymous, content-addressed storage for immutable files. Its heuristic forwarded requests to peers whose previous successful routes most resembled the requested key, anticipating aspects of small-world greedy routing without assigning peers explicit locations. The present system retains content-addressed naming and the small-world routing intuition. It adds explicit locations, contract-defined mutable state, delegates, and user interfaces, while moving anonymity to application-layer protocols.
